\section{Prompt}~\label{appx:prompt}

\begin{tcolorbox}[breakable, title=Prompt for feedback manager (Waste Sorting), title filled=false, colback=black!5!white, colframe=black!75!black, left=2pt, right=2pt, top=2pt, bottom=2pt]
\textbf{System}: You convert one symbolic waste-sorting fact into exactly one concise Korean yes/no question. You must preserve the predicate meaning exactly. Do not change the predicate into another property. Ask only the final question in Korean.
\\
\\
\textbf{Message}: Instruction: Convert the given symbolic fact into a polite and unambiguous Korean yes/no question.\\
\\
Important rules:\\
1. First identify the predicate name in the fact.\\
2. The question must ask about that predicate only.\\
3. Do not infer or replace it with another predicate.\\
4. If the fact is detected(W), ask whether the indicated waste item has been detected.\\
5. If the fact is plastic(W), can(W), paper(W), or general(W), ask about the indicated waste item's type.\\
6. If the fact is holding(R, W), ask whether the robot is holding the indicated waste item.\\
7. If the fact is handempty(R), ask whether the robot's hand is empty.\\
8. If the fact is in\_bin(W, B), ask whether the indicated waste item is in the indicated bin.\\
9. The respondent may not recognize identifiers such as waste1, waste2, robot1, bin\_01, or bin\_02.\\
10. Assume the respondent sees only a tablet display where the relevant object, robot, or bin is visually indicated.\\
11. Output only the question.\\
\\
Examples:\\
Fact: detected(waste1)\\
Question: 화면에 표시된 폐기물이 감지된 것이 맞나요?\\
\\
Fact: plastic(waste1)\\
Question: 화면에 표시된 폐기물이 플라스틱인 것이 맞나요?\\
\\
Fact: holding(robot1,waste2)\\
Question: 현재 로봇이 폐기물을 들고 있는 것이 맞나요?\\
\\
Fact: in\_bin(waste1,bin\_01)\\
Question: 화면에 표시된 폐기물이 표시된 수거함 안에 있는 것이 맞나요?\\
\\
Action: \cp{\{selected action\}}\\
Fact: \cp{\{selected hypothesis\}}\\
\\
Predicate meanings:\\
- detected(W): waste item W has been detected.\\
- plastic(W): waste item W is plastic.\\
- can(W): waste item W is a can.\\
- paper(W): waste item W is paper.\\
- general(W): waste item W is general waste.\\
- holding(R, W): robot R is holding waste item W.\\
- handempty(R): robot R's hand is empty.\\
- in\_bin(W, B): waste item W is in bin B.\\
\\
Question:
\\
\\
\textbf{Answer: } 화면에 표시된 폐기물이 감지된 것이 맞나요?

\end{tcolorbox}



\begin{tcolorbox}[breakable, title=Prompt for feedback manager (Tomato Harvesting), title filled=false, colback=black!5!white, colframe=black!75!black, left=2pt, right=2pt, top=2pt, bottom=2pt]
\textbf{System}: You convert one symbolic tomato-harvesting fact into exactly one concise Korean yes/no question. You must preserve the predicate meaning exactly. Do not change the predicate into another property. Ask only the final question in Korean.
\\
\\
\textbf{Message}: Instruction: Convert the given symbolic fact into a polite and unambiguous Korean yes/no question.\\
\\
Important rules:\\
1. First identify the predicate name in the fact.\\
2. The question must ask about that predicate only.\\
3. Do not infer or replace it with another predicate.\\
4. If the fact is located(R, L), ask whether the robot is at the indicated location.\\
5. If the fact is ripe(T), unripe(T), or rotten(T), ask about the indicated tomato's state.\\
6. If the fact is at(T, S), ask whether the indicated tomato is attached to or located at the indicated stem/location.\\
7. The respondent may not recognize identifiers such as stem\_01, stem\_02, tomato1, or tomato2.\\
8. Assume the respondent sees only a tablet display where the relevant object or location is visually indicated.\\
9. Output only the question.\\
\\
Examples:\\
Fact: located(brl\_robot,stem\_01)\\
Question: 지금 로봇이 화면에 표시된 위치에 있는 것이 맞나요?\\
\\
Fact: ripe(tomato1)\\
Question: 화면에 표시된 토마토가 익은 것이 맞나요?\\
\\
Fact: rotten(tomato2)\\
Question: 화면에 표시된 토마토가 상한 것이 맞나요?\\
\\
Fact: at(tomato1,stem\_??)\\
Question: 현재 줄기에 첫번째 토마토가 달려있나요?\\
\\
Action: \cp{\{selected action\}}\\
Fact: \cp{\{selected hypothesis\}}\\
\\
Predicate meanings:\\
- observed(T): tomato T has been observed.\\
- scanned(T): tomato T has been scanned.\\
- ripe(T): tomato T is ripe.\\
- unripe(T): tomato T is unripe.\\
- rotten(T): tomato T is rotten.\\
- holded(T, R): robot R has held tomato T.\\
- loaded(T, R): tomato T has been loaded by robot R.\\
- discarded(T): tomato T has been discarded.\\
- at(T, S): tomato T is at stem/location S.\\
- handempty(R): robot R's hand is empty.\\
- holding(R, T): robot R is holding tomato T.\\
- located(R, L): robot R is located at location L.\\
\\
Question:
\\
\\
\textbf{Answer: } 화면에 표시된 토마토가 익은 것이 맞나요?

\end{tcolorbox}